\documentclass[12pt, spanish]{article}
\usepackage[utf8]{inputenc}
\usepackage[T1]{fontenc}
\usepackage[spanish]{babel}
\usepackage{amsmath, amssymb}
\usepackage{geometry}
\geometry{left=2.5cm, right=2.5cm, top=2.5cm, bottom=2.5cm}
\usepackage{graphicx}
\usepackage{float}
\usepackage{array}
\usepackage{booktabs}
\usepackage{enumitem}
\usepackage{xcolor}
\usepackage{hyperref}
\hypersetup{
    colorlinks=true,
    linkcolor=blue,
    urlcolor=blue,
}

\title{Informe de Producción\\Shadow-Score Académico (v2)}
\author{Equipo de desarrollo}
\date{26 de marzo de 2026}

\begin{document}

\maketitle

\begin{abstract}
El proyecto Shadow-Score Académico consiste en una aplicación web interactiva que permite a estudiantes universitarios reflexionar sobre la distribución de su tiempo entre cargas domésticas, laborales y académicas, y visualizar posibles afectaciones en su rendimiento académico desde una perspectiva de género. Esta versión 2 incorpora un modelo matemático multifactorial basado en la retroalimentación de docentes de investigación. La solución se construye con Streamlit y se describe en detalle: arquitectura, modelo ajustado, interfaz, componentes y cronograma.
\end{abstract}

\section{Introducción}

El proyecto Shadow-Score Académico consiste en una \textbf{aplicación web interactiva} que permite a estudiantes universitarios reflexionar sobre la distribución de su tiempo entre cargas domésticas, laborales y académicas, y visualizar posibles afectaciones en su rendimiento académico desde una perspectiva de género. La herramienta se construye con base en los fundamentos investigativos desarrollados previamente y responde a la retroalimentación recibida, ampliando su enfoque hacia la \textbf{autorreflexión} y considerando la \textbf{multifactorialidad} del fenómeno.

El presente informe de producción describe la solución técnica a programar: su arquitectura, modelo matemático ajustado, interfaz de usuario, componentes, tecnologías y cronograma de desarrollo. El resultado final será una aplicación desplegada en Streamlit Cloud, accesible públicamente.

\section{Objetivos de producción}

\subsection{Objetivo general}
Desarrollar una aplicación web interactiva que integre un modelo multifactorial de cálculo de fatiga y potencial afectación académica, basado en variables de carga doméstica, laboral, académica y contextuales, con un enfoque de autorreflexión y sensibilización sobre las brechas de género.

\subsection{Objetivos específicos de producción}
\begin{enumerate}
    \item Implementar un modelo matemático que considere carga doméstica, carga laboral, composición del hogar, estrato, género y percepción subjetiva, para generar indicadores de fatiga, tiempo efectivo de estudio y un Shadow-Score cualitativo.
    \item Diseñar una interfaz de usuario intuitiva que recoja las variables definidas en la investigación y presente resultados mediante gráficos interactivos, métricas y elementos reflexivos.
    \item Simular escenarios de corresponsabilidad y apoyos institucionales (reducción de carga doméstica, acceso a servicios de cuidado) que permitan comparar situaciones actuales y mejoradas.
    \item Incluir información contextual (datos DANE, referencias a ODS) y recomendaciones de apoyo institucional, fomentando la reflexión y la toma de conciencia.
    \item Garantizar la escalabilidad, mantenibilidad y accesibilidad de la aplicación, utilizando buenas prácticas de desarrollo.
\end{enumerate}

\section{Arquitectura técnica}

La aplicación se construirá en \textbf{Python} utilizando \textbf{Streamlit} como framework de frontend.

\subsection{Diagrama de componentes}
\begin{verbatim}
[Usuario] → (Streamlit UI) → (Validación de datos) → (Modelo matemático)
                                                           ↓
[Visualización y reflexión] ← (Resultados) ← (Simulación de escenarios)
\end{verbatim}

\subsection{Tecnologías y librerías}
\begin{table}[H]
\centering
\begin{tabular}{ll}
\toprule
\textbf{Tecnología} & \textbf{Uso} \\
\midrule
Python 3.10+ & Lenguaje base \\
Streamlit & Interfaz de usuario, gestión de estado, despliegue \\
Pandas & Manejo de datos (almacenamiento temporal) \\
NumPy & Operaciones matemáticas y vectorización \\
Plotly & Gráficos interactivos (barras, radar, medidores) \\
Streamlit-extras & Componentes adicionales (gauge, métricas) \\
Git / GitHub & Control de versiones \\
Streamlit Cloud & Despliegue continuo \\
\bottomrule
\end{tabular}
\end{table}

\subsection{Estructura de carpetas}
\begin{verbatim}
shadow-score/
+-- .streamlit/
|   +-- config.toml
+-- src/
|   +-- __init__.py
|   +-- modelo.py
|   +-- utils.py
|   +-- graficos.py
+-- assets/
|   +-- logo.png
|   +-- estilo.css
+-- app.py
+-- requirements.txt
+-- README.md
+-- .gitignore
\end{verbatim}

\section{Modelo matemático ajustado (multifactorial)}

El modelo se ha ampliado para reflejar la multifactorialidad. Se basa en parámetros derivados de la literatura (DANE ENUT, estudios colombianos, ODS).

\subsection{Variables de entrada}
\begin{table}[H]
\centering
\footnotesize
\setlength{\tabcolsep}{3pt}
\begin{tabular}{>{\raggedright\arraybackslash}p{0.15\textwidth} >{\raggedright\arraybackslash}p{0.15\textwidth} >{\raggedright\arraybackslash}p{0.30\textwidth} >{\raggedright\arraybackslash}p{0.28\textwidth}}
\toprule
\textbf{Variable} & \textbf{Tipo} & \textbf{Rango / Opciones} & \textbf{Notas} \\
\midrule
Género & Categórica & Femenino / Masculino / Otro & Ajuste de parámetros \\
Estrato & Numérico entero & 1–6 & Factor de ponderación \\
Composición del hogar & Categórica & Vive solo/a, con familia, con pareja, residencia universitaria, otros & Modula apoyo \\
Dependientes & Numérico entero & 0, 1, 2, 3+ & Personas a cargo \\
Horas domésticas & Numérico (float) & 0–168 & Desglose opcional \\
Trabajo remunerado & Sí/No + horas & 0–168 & \\
Horas de estudio & Numérico & 0–168 & \\
Promedio actual (opcional) & Numérico & 0–5.0 & Para comparación \\
\bottomrule
\end{tabular}
\end{table}

\subsection{Cálculo de carga total y fatiga cognitiva}
\[
\text{carga\_total} = \text{horas\_domésticas} + \text{horas\_trabajo}
\]
\[
\text{fatiga} = \sigma\left(\frac{\text{carga\_total} - \mu}{\kappa}\right) \cdot \alpha_{\text{género}} \cdot \beta_{\text{estrato}} \cdot \gamma_{\text{hogar}}
\]
donde \(\sigma(x)=\frac{1}{1+e^{-x}}\), y los parámetros \(\mu, \kappa, \alpha, \beta, \gamma\) están precalibrados con fuentes secundarias.

\subsection{Horas de estudio efectivas}
\[
\text{horas\_efectivas} = \text{horas\_estudio} \times (1 - \text{fatiga})
\]

\subsection{Estimación del promedio (PPA)}
\[
\text{PPA\_estimado} = \beta_0 + \beta_1 \cdot \text{horas\_efectivas} + \beta_2 \cdot \text{estrato} + \beta_3 \cdot \text{género} + \beta_4 \cdot \text{dependientes} + \beta_5 \cdot (\text{género} \times \text{fatiga})
\]

\subsection{Shadow-Score}
\[
\text{Shadow-Score} = \min\left(100,\; \frac{\text{carga\_total}}{168} \times 100 \times \text{fatiga}\right)
\]
Interpretación: 0–20 (baja penalización), 21–40 (moderada), 41–60 (significativa), 61–80 (alta), 81–100 (extrema).

\subsection{Escenario de corresponsabilidad}
Se simula una reducción en la carga doméstica (por defecto 40\%, ajustable) y opcionalmente un apoyo institucional que reduce la fatiga. Se comparan indicadores lado a lado.

\section{Flujo de usuario y funcionalidades}

\begin{enumerate}
    \item \textbf{Página de inicio}: presentación, ODS 5, datos DANE.
    \item \textbf{Formulario de entrada}: datos sociodemográficos, carga de trabajo, académicos y percepción.
    \item \textbf{Cálculo y visualización}: Shadow-Score, fatiga (gauge), horas efectivas, promedio estimado, gráfico de radar.
    \item \textbf{Simulación}: escenario actual vs. mejorado.
    \item \textbf{Información adicional}: metodología, fuentes, recursos institucionales.
\end{enumerate}

\section{Plan de desarrollo (cronograma)}

\begin{table}[H]
\centering
\small
\begin{tabular}{clp{8cm}}
\toprule
\textbf{Semana} & \textbf{Fase} & \textbf{Actividades} \\
\midrule
1 & Configuración y modelo & Repositorio, dependencias, funciones básicas, pruebas unitarias. \\
2 & Interfaz básica & Formulario simple, validaciones, integración del modelo. \\
3 & Visualizaciones & Gráficos Plotly (medidor, radar, barras). \\
4 & Simulación & Lógica de reducción de carga y apoyo institucional, vista comparativa. \\
5 & Elementos reflexivos & Expanders con datos DANE, referencias, recomendaciones. \\
6 & Pruebas de usuario & Feedback, correcciones. \\
7 & Optimización & Caché, tiempos de respuesta. \\
8 & Despliegue & Streamlit Cloud, pruebas multi-dispositivo. \\
9 & Documentación & README, video demostrativo. \\
\bottomrule
\end{tabular}
\end{table}

\section{Consideraciones de calidad}
\begin{itemize}
    \item Código limpio (PEP 8, docstrings, typing).
    \item Pruebas unitarias para funciones críticas.
    \item Manejo de errores y validaciones.
    \item Accesibilidad y contraste.
    \item Privacidad: no almacenar datos personales sin consentimiento.
\end{itemize}

\section{Ampliaciones futuras}
\begin{itemize}
    \item API REST con FastAPI.
    \item Almacenamiento anónimo de datos para recalibrar el modelo.
    \item Recomendaciones personalizadas.
    \item Dashboard institucional.
\end{itemize}

\section{Conclusión}
La versión 2 de Shadow-Score Académico presenta un modelo multifactorial robusto, basado en evidencia y retroalimentación, que permite una reflexión más realista sobre la carga doméstica y el rendimiento académico. La implementación con Streamlit garantiza una herramienta accesible, interactiva y de fácil despliegue.

\end{document}